\chapter[Religiöser Unterricht]{Religiöser Unterricht der Jugend}
\quellseite{7}
Daß die Jugend unseres Landes besser in der Religion unterrichtet werden muß, als dies in der einen kurzen Stunde am Sonntag in der Sonntagsschule geschehen kann, wird nachgerade fast in allen amerikanisch-kirchlichen Kreisen zugestanden. Wir haben schon viele dahinlautende Aussprüche mitgetheilt und könnten in jeder Nummer neue bringen. Ein bekannter Prediger Bostons klagte vor nicht langer Zeit auf einer dortigen Conferenz, daß es Gemeinden gebe, die wohl \$2000.00 im Jahr einem aus vier Personen bestehenden Singchor bezahlen, der dreimal im Gottesdienst zu singen habe, aber nur \$250.00 ausgäben, um 500 ihnen anvertraute Kinder und junge Leute religiös zu erziehen. Dr. Rufus Miller, der Secretär der Reformirten Sonntagsschulbehörde, also gerade ein Vertreter der Sonntagsschule, schreibt in der Christian World: „Es leuchtet ein, daß Angesichts der veränderten Verhältnisse der heutigen Gesellschaft irgend etwas gethan werden muß, damit die jungen Leute einen planmäßigeren und klareren Unterricht in der biblischen Geschichte und in den christlichen Lehren erhalten.“ Der Episcopälerprediger Dr. Montague Geer in New York schrieb vor einigen Monaten in der New York Sun einen Aufsehen erregenden Artikel. Er möchte am liebsten unser jetziges öffentliches Schulsystem abgeschafft sehen, weil wir in unserm Lande durch die religionslose Schule „ein Geschlecht von jungen Heiden erziehen“. Er betont, daß „das Recht und die Pflicht der Kindererziehung den Eltern gehört, welche es auch der Kirche und den Privatlehrern überlassen können. Auf diese Weise wird den Kirchengemeinden freie Hand gelassen im Unterrichtswesen, und zahlreiche Kirchen- und Privatschulen werden entstehen, die der Jugend Religion und Patriotismus beibringen werden. Das ist der richtige Standpunkt im Schulwesen. Heiden mögen ihre Kinder erziehen im Heidenthum, Juden die ihrigen im Judenthum, Protestanten die ihrigen im Protestantismus, Katholiken die ihrigen im Katholicismus“. Bei einer Methodistenversammlung in Cincinnati sagte der Redacteur des Christian Advocate: „Die Römischen, die Episcopalen und die Lutheraner haben längst eingesehen, wie wichtig es ist, die Kinder in den Grundlehren des christlichen Glaubens gründlich zu unterrichten, und die Methodisten verlieren jedes Jahr Tausende von Kindern, weil sie nicht früher eingesehen haben, wie wichtig diese Sache ist.“ Ganz ähnlich äußerte sich im Advance, einem Baptistenblatte, ein Schreiber über die Kirchen, welche ihre Kinder im Katechismus unterrichten, und sagt: „Sie erziehen ihre Kinder im Katechismus, in der Lehre des Wortes Gottes; wenn diese dann älter werden, so fallen sie nicht ab von ihrer Kirche. Es ist das eine so vernünftige Idee, daß man sich wundern muß, weshalb wohl so viele Kirchen von dieser Weise abgewichen sind.“ Richtig bemerkt zur letzteren Aussprache die „Lutherische Kirchenzeitung“, daß dies nicht nur eine vernünftige Idee, sondern ausdrücklicher Befehl Gottes ist, die Kinder aufzuziehen in der Zucht und Vermahnung zum Herrn, Eph. 6, 4. Freilich, an die einzig richtige und beste Lösung der Frage, an die Errichtung von Gemeindeschulen, will man nicht heran. Der Lutheran des Generalconcils, der die Sache der besseren religiösen Erziehung der Kinder eine „beunruhigende Frage“ nennt, bezeichnet immer wieder dies als Lösung der Schwierigkeit, daß die einzelnen Kirchen mit den Schulbehörden Vereinbarungen treffen sollten, daß mehrere Schulstunden während der Woche frei bleiben, in denen dann jede Kirche ihre Kinder in der Religion unterrichten lassen kann. Auf diese Weise würde der größte Mangel der Staatsschulen beseitigt, ohne sie doch aufzugeben. Der Lutheran freut sich darüber, daß sich ein so angesehener Mann wie der bekannte Präsident Schurmann von der Cornell-Universität kürzlich ganz in demselben Sinne ausgesprochen habe. Aber dieser Plan wird schwerlich in absehbarer Zeit allgemein durchgeführt werden und ist vor allem noch nicht genügend, denn die Kinder sollen nicht nur zwei oder drei Stunden in der Woche in der Religion unterrichtet werden, sondern täglich, und sollen auch in den übrigen Schulstunden unter christlichem Einfluß stehen. L. F.
